% !!! GENERATED FILE - DO NOT EDIT !!!
% Produced from paper/sections/10_appendix_d.html by paper/assemble_arxiv.py
% (rules in arxiv/_CONVERSION_SPEC.md). Edit the .html source + meta.py, then
% rerun `python paper/build.py` (or `python paper/assemble_arxiv.py`).

\section{Appendix D. Extended out-of-loop results}\label{sec:appendixd}

The body (\cref{sec:payoff}) carries the two headline live results --- the write-admission
gate (J~=~10) and the race arbiter (J~=~6). This appendix holds three further
results that \emph{strengthen} the existence claim without being the headline: a richer witness and
the first measurement of what one propagated phantom costs (D.1); a frontier test that puts a frontier
model in the verifier's own seat against a ladder of non-forgeable rungs (D.2); and the
compounding payoff across a branching fleet (D.3). All three read the same kind of byte the agent
cannot author, and each carries its evidence-strength label in the same breath as its number.

\subsection{D.1 A richer witness, and the first measurement of what propagation costs}\label{sec:appendixd-richerwitness}

tau2's witness is a single bit, so a reader might ask whether the result is an artifact of that
coarse signal. We re-aimed the \emph{identical} gate at a second benchmark with a much richer witness:
\textbf{Agent-Diff} (a write-heavy enterprise-API benchmark over Slack, Box, Linear, and Google
Calendar), whose environment computes a full \emph{structured state-diff} from database snapshots ---
the exact rows inserted, updated, deleted --- and grades against a multi-assertion gold
specification. Driving an agent live through Gemini on the 45-task test split, the
\textbf{detection result holds across the capability step a second time}: flash over-claims on
23 of its 35 confident writes (66\%), pro on
25 of 31 (81\% --- \emph{more}, not less), and the gate blocks
every one. The over-claims are partly \emph{model-specific} (14 tasks over-claim on both models, the
rest on only one), so the gate catches the \emph{deployed} model's blind spots, not a fixed set
of impossible tasks. Detection generalizes off the coarse \code{db\_match} bit to a structured
diff and across two tiers.

Agent-Diff also lets us measure the thing J names but does not itself quantify: \textbf{what does
letting a phantom write propagate actually cost a downstream agent?} We build the peer~B
handoff explicitly. For each blocked over-claim we run a second agent twice from the \emph{same}
fresh, correct environment, changing only what B is told it inherited: under \textbf{believe}
B receives A's claim verbatim; under \textbf{adjudicate} B receives the gate's
env-verified correction (``the prior write did not land --- re-verify before acting'').
The measured quantity is

\textbf{$\Delta$B = success(B~|~adjudicate) $-$ success(B~|~believe)}

over the over-claim slice. On flash, \textbf{$\Delta$B~=~+4}
(four flips, zero reverses): telling the peer the truth turned four phantom-induced failures into
successes and never made one worse. On pro, \textbf{$\Delta$B~=~0}
(two flips cancelled by two reverses; one reverse traced to a real gate bug --- a blanket
``treat the records as unchanged'' over-corrected a \emph{partial} over-claim in which one
of two writes had landed. We fixed the gate for future runs but report the as-run number).

We lead with J, not $\Delta$B, because they are a product and a lower bound on one of its terms, not
two estimates of one thing. J counts the lies stopped at the root; $\Delta$B measures the cost of
letting \emph{one} lie propagate \emph{one} hop to the most \emph{forgiving} consumer imaginable
--- an LLM peer that re-runs the task and can sometimes heal the missing write for free. So
$\Delta$B is a pessimistic floor: the gate's real value is \textbf{blocks $\times$ depth $\times$ per-hop
damage}, and $\Delta$B pins depth and consumer to their minimum. A consumer that does \emph{not}
re-check (a reward labeler, an auto-merge, a dashboard) recovers nothing. The \emph{depth} term is
what D.3 measures.

\subsection{D.2 The frontier test: a ladder of rungs a frontier model cannot all climb}\label{sec:appendixd-frontierfloor}

The cross-model results answer ``does a \emph{stronger agent} over-claim less'' ---
it does not. The sharper version of the objection is that \textbf{the latest model is smart enough that
you no longer need a referee} --- either it stops over-claiming, or, as a judge, it sees
through forgeries an older checker missed. We test the decisive second half: the referee here has two
implementations, the deterministic floor that reads ground truth or a model asked to \emph{judge} the
same claim, so we put a frontier model in the judge seat and raced it against the floor. The model is
\textbf{Fable~5} (the newest at the time of writing), reached the way an operator would
--- straight through the Claude~Code command line (\code{claude -p}), the model
named only in an environment variable, \emph{not a line of model-specific code} in the harness.

The test materializes nine artifacts as \emph{real, freshly authored} git commits --- seven
forgeries plus two honest controls. Four forgeries touch the \emph{wrong} files (an empty commit
claiming the work; a keyword-stuffed subject; deleting the test that would have caught the gap; a long
``production-ready'' narration over a commit that did none of it). The other three are the
adversary that matters --- commits that touch the \textbf{exact deliverable files the plan
names} but ship nothing behind them (empty stubs; code that does not import; a lookalike
filename). Both channels see the \emph{same} forgery-resistant evidence; the only asymmetry, stated up
front so the fight is fair, is that the model judge \emph{also} reads the agent's narration. The
floor is a \textbf{ladder} of three rungs: the commit \emph{subject} (forgeable, the author
types it); \emph{which files the commit touched} (a commit cannot fake that, but it certifies
presence, not behavior --- a stub passes); and the OS \emph{running the test and reading the exit
code} (a \code{0} the agent cannot forge, because it did not run the process).

\begin{figure*}[t]
  \centering
  \includegraphics[width=\textwidth]{forge_frontier_floor}
  \caption{\cref{fig:frontierfloor}. \textbf{The floor is a ladder, and its top rung still
catches a frontier model's best forgery.} \textbf{Left:} on the
8 not-shipped artifacts, the grep-the-subject rung is fooled
7/8 and the presence rung
2/8 (path-touching stubs), but the behavior rung ---
the OS runs the test and reads the exit code --- holds at
0/8. \textbf{Right:} on the one honest landing
the same Fable judge abstains given narration or files-touched
(0/1) but confirms it given the OS test
result (1/1).}
  \label{fig:frontierfloor}
\end{figure*}

The result is not ``the floor is perfect.'' The \emph{presence} rung --- what most
people picture by ``check that the files landed'' --- is \textbf{fooled
2 of 8} by the path-touching stubs, the honest
ceiling of any file-counting check. What sits \emph{above} it is the point: the \textbf{behavior rung
holds at 0 of 8} and still confirms the one real
landing, because it has the OS \emph{run the code} and reads an exit code the agent cannot author (the
forgeable subject rung is fooled on 7 of 8). The lesson is
not ``determinism wins'' but \textbf{which non-forgeable byte you condition belief on, and that
the rung you must stand on rises with the forger's sophistication.}

The judge-side result is the more surprising one: the learned judge is \emph{exactly as good as the
evidence the kernel computes for it}. Given only narration, Fable abstains on the honest landing
(0/1); given only which files were touched it
\emph{still} abstains (0/1) --- it
will not certify presence as completion. But given the OS test result --- the byte the behavior rung
reads --- the \emph{same model} confirms the work
(1/1) and disagrees with every forgery
(0/8). Left to its own reading the judge yields no value
--- gamed if forced to answer (35.2\% false-accepts elsewhere), inert if allowed to abstain.
\textbf{The floor's value is that it \emph{computes} the non-forgeable witness --- including
running the code --- for free, every time, lifting both the deterministic verdict and the learned
one.} The stronger the model, the more the value concentrates in the byte it cannot author.
(Full grid: docs/272.)

\subsection{D.3 The payoff compounds: one block at the root keeps the whole downstream fleet clean}\label{sec:appendixd-cascade}

$\Delta$B priced a \emph{single} phantom propagating \emph{one} hop. A real fleet is neither
single nor shallow --- a phantom is inherited again by the next agent, and a fleet
\emph{branches} --- so the true quantity the gate protects is the whole set of downstream agents a
root corruption reaches. We measured that set live, down a chain and across a fan-out tree, on the same
tau2 airline witness with two arms (\emph{believe}: each agent inherits the corrupt state;
\emph{adjudicate}: the gate blocked the root, so each inherits the gold state). The root corruption is
one reservation~R wrongly cancelled (injected; every \emph{downstream} node runs live), and the
witness is exact: \textbf{is R still cancelled when this agent finishes?}

\begin{figure*}[t]
  \centering
  \includegraphics[width=\textwidth]{cascade_depth_live}
  \caption{\cref{fig:cascade}. The cascade, run live, down a chain. Under
\textbf{believe} all D$-$1 downstream agents inherit the cancelled reservation and are
left corrupt (rising 1--2--3 with depth); under \textbf{adjudicate} the one root
block keeps every node on the gold database (flat at zero). The believe agents \emph{did not heal
it} --- a corrupt \emph{state} is not second-guessed away the way a narrated claim is,
which is why $\Delta$B was near zero to a self-recovering peer.}
  \label{fig:cascade}
\end{figure*}

Down a single chain the result is linear: under \emph{believe} a depth-D chain leaves all
\textbf{D$-$1} downstream nodes corrupt, under \emph{adjudicate} zero (the gate prevents
1, 2, and 3 corrupt nodes at depths 2, 3, 4). The believe agents are
instructive: each, handed a cancelled reservation, simply could not do its task and gave up --- none
restored R. So a capable peer re-checks and heals a \emph{narrated} phantom (hence
$\Delta$B~$\approx$~0) but cannot re-check away a corrupt \emph{starting world}; the state
handoff is the one that compounds.

\begin{figure*}[t]
  \centering
  \includegraphics[width=\textwidth]{cascade_fanout_live}
  \caption{\cref{fig:cascadetree}. The same cascade across a fan-out tree
(F~=~2). \textbf{Left:} the chain's linear D$-$1 payoff (1, 2) against
the tree's F$^{D}$ (4, 8). Branching makes
the payoff super-linear in depth. \textbf{Right:} the depth-3 tree --- under
\emph{believe} all 8 leaves are corrupt; under \emph{adjudicate} the one block at the root keeps
all 8 clean.}
  \label{fig:cascadetree}
\end{figure*}

A chain understates a fleet, because fleets branch. A depth-D, fan-out-F dependency \emph{tree} has
F$^{D}$ leaves, all downstream of the root corruption. Under \emph{believe} every leaf is corrupt
(\textbf{4 of 4} at depth~2, \textbf{8 of 8} at depth~3, F~=~2);
under \emph{adjudicate} the single root block keeps every leaf clean. So the prevented-corruption
payoff is \textbf{F$^{D}$} --- 4 at depth~2,
8 at depth~3 --- against the chain's linear D$-$1. F$^{D}$
here is \textbf{breadth} (F$^{D}$ agents each \emph{blocked} by the one shared poison,
not one corrupt value mutating into many): a poisoned shared resource fails every agent that depends on
it, and that set grows F$^{D}$ as the fleet branches. This is the \emph{depth} term $\Delta$B
pinned to one, now measured live. (Caveats: the root is injected, not waited for, and each node
re-inherits the same root state rather than chaining onto the previous node's output; it is an
existence result for the compounding \emph{shape} --- one model, F~=~2, depths 2--4
--- not a calibrated constant. The fleet-of-one floor holds: at depth~0 the payoff is zero.)
